\Subhead{4.1.1}{Принцип работы, устройство и параметры: гелий-неоновый лазер, лазер на рубине, полупроводниковый лазер}

Лазер преобразует энергию накачки в когерентное поле через вынужденное излучение и обратную связь. Для двух уровней с населённостями $N_2,N_1$ усиление определяется разностью $N_2-(g_2/g_1)N_1$, где $g_j$ — вырождения. Тепловое равновесие инверсии не даёт. В чистой замкнутой двухуровневой системе резонансная накачка вызывает также обратные переходы и не создаёт стационарной положительной инверсии: нужны дополнительные уровни или неравновесное распределение носителей.

В простом четырёхуровневом лазере быстро опустошаемый нижний рабочий уровень упрощает создание инверсии. Для числа фотонов $S$ и эффективной инверсии $N$
\begin{equation}
 \dot S=(GN-\kappa)S+R_{\rm sp},\qquad
 \dot N=R_p-N/\tau-GNS.
\end{equation}
$\kappa$ — скорость потери фотонов, $R_p$ — накачка, $R_{\rm sp}$ — спонтанный источник выбранной моды. При малом поле порог $GN=\kappa$; выше порога инверсия приблизительно зажимается, дополнительная накачка увеличивает поле. Подробный вывод и модальная конкуренция — в вопросах~\Qref{1.10.1}--\Qref{1.10.2}.

\begin{center}\small
\begin{tabular}{p{.16\linewidth}p{.30\linewidth}p{.42\linewidth}}
\toprule
Лазер & Создание инверсии & Физическая особенность\\\midrule
He--Ne & Газовый разряд возбуждает He; резонансные столкновения передают энергию Ne & Рабочее излучение Ne; известная красная линия около 632.8 нм. Низкий коэффициент усиления требует малых потерь резонатора.\\
Рубиновый & Оптическая накачка Cr$^{3+}$ в Al$_2$O$_3$; накопление на метастабильном уровне & Трёхуровневая схема: нижний рабочий уровень — основной, высокий порог; линия около 694.3 нм, удобен импульсный режим.\\
Полупро\-воднико\-вый & Инжекция электронов и дырок в активный слой гетероструктуры & Разность квазиуровней Ферми обеспечивает усиление; оптическое и электронное ограничение снижают порог. Длина волны задаётся зонной структурой.\\\bottomrule
\end{tabular}\end{center}
Для инжекционного лазера выше порога простая оценка выходной мощности $P_{\rm out}\simeq\eta_d(h\nu/e)(I-I_{\rm th})$. Здесь $\eta_d$ — внешняя дифференциальная квантовая эффективность выбранного выхода, а не полный электрический КПД. Параметры для сравнения: порог, мощность, дифференциальная эффективность, ширина спектра, пространственная мода, расходимость, быстродействие и тепловая чувствительность. Конкретные численные характеристики зависят от устройства, а не только от названия активной среды.
